% GENERATED FILE. Edit canonical lessons, then run npm run generate:books.
% canonical-source-hash: fnv1a64:0c55f420f455ad5a
% canonical-lessons: MR-C467-yatna, MR-C467-tadphadne, MR-C467-kalvalne, MR-C467-jagne-live, MR-C467-dagmagne

\chapter{Bearing It}
\label{ch:mr-bearing-it}
\hlchaptermodality{467}

\begin{chapteropening}
\textbf{By the end of this chapter:} \emph{I can talk about torment, writhing, pity and wavering.}

Complete the last lesson of chapter 467: I can talk about torment, writhing, pity and wavering.
\end{chapteropening}

\section[\texorpdfstring{\mr{यातना} (yātnā)}{यातना (yātnā)}]{\mr{यातना} (yātnā) — torment; agony}

\label{lesson:MR-C467-yatna}



\begin{quote}
Before the new one: say the Marathi for to rest a while, then the Marathi for to brighten up.
\end{quote}

\subsection*{You'll want to know: \mr{यातना}}
\textbf{\mr{यातना}} — \emph{yātnā} — \textquotedblleft{}torment; agony\textquotedblright{}.

\begin{quote}
\mr{आजारपणात} \mr{आजोबांनी} \mr{खूप} \mr{यातना} \mr{सहन} \mr{केल्या}.
\end{quote}

\begin{quote}
\emph{ājārpaṇāt ājobānnī khūp yātnā sahan kelyā.}
\end{quote}

\begin{quote}
Grandfather bore great agony during his illness.
\end{quote}

Feminine.

Often plural: \mr{यातना} \mr{सहन} \mr{करणे}, to bear agony.

\begin{cousinweb}
\textbf{\mr{यातना}} (\emph{yātnā}) is a tatsama, taken from Sanskrit \emph{yātanā}, \textquotedblleft{}torment\textquotedblright{} (Monier-Williams).
\end{cousinweb}

\begin{grammarlens}[title={What you've built}]
One more word for describing what you lived and felt.
\end{grammarlens}

\subsection*{Your turn}
\begin{itemize}
\raggedright
  \item \emph{Say it:} \emph{yātnā}
  \item \emph{Say it:} \emph{yātnā}, once more
  \item \emph{Say it:} \emph{ājārpaṇāt ājobānnī khūp yātnā sahan kelyā.} — Grandfather bore great agony during his illness.
  \item \emph{Recall:} say the Marathi for to rest a while, then the Marathi for to brighten up, then say \emph{yātnā} again
\end{itemize}

\begin{tcolorbox}[breakable,colback=teal!4,colframe=teal!35!black,title={Before you move on}]
What does \mr{यातना} mean? (\textquotedblleft{}Torment; agony\textquotedblright{}.) Say it once more.
\end{tcolorbox}
\section[\texorpdfstring{\mr{तडफडणे} (taḍphaḍṇe)}{तडफडणे (taḍphaḍṇe)}]{\mr{तडफडणे} (taḍphaḍṇe) — to writhe}

\label{lesson:MR-C467-tadphadne}



\begin{quote}
Before the new one: say the Marathi for to brighten up, then the Marathi for torment.
\end{quote}

\subsection*{You'll want to know: \mr{तडफडणे}}
\textbf{\mr{तडफडणे}} — \emph{taḍphaḍṇe} — \textquotedblleft{}to writhe\textquotedblright{}.

\begin{quote}
\mr{पाण्याबाहेर} \mr{काढलेला} \mr{मासा} \mr{तडफडत} \mr{होता}.
\end{quote}

\begin{quote}
\emph{pāṇyābāher kāḍhlelā māsā taḍphaḍat hotā.}
\end{quote}

\begin{quote}
The fish taken out of the water was writhing.
\end{quote}

Past (he …): \mr{तो} \mr{तडफडला}, \emph{to taḍphaḍlā}.

Of pain, or of fierce impatience: \mr{भेटायला} \mr{तडफडणे}.

\begin{cousinweb}
\textbf{\mr{तडफडणे}} (\emph{taḍphaḍṇe}) is inherited: CDIAL lists Marathi \emph{taḍphaḍṇẽ}, \textquotedblleft{}to toss about\textquotedblright{}, under the expressive \emph{*taḍapphaḍ} (CDIAL 5634).
\end{cousinweb}

\begin{grammarlens}[title={What you've built}]
One more word for describing what you lived and felt.
\end{grammarlens}

\subsection*{Your turn}
\begin{itemize}
\raggedright
  \item \emph{Say it:} \emph{taḍphaḍṇe}
  \item \emph{Say it:} \emph{taḍphaḍṇe}, once more
  \item \emph{Say it:} \emph{pāṇyābāher kāḍhlelā māsā taḍphaḍat hotā.} — The fish taken out of the water was writhing.
  \item \emph{Recall:} say the Marathi for to brighten up, then the Marathi for torment, then say \emph{taḍphaḍṇe} again
\end{itemize}

\begin{tcolorbox}[breakable,colback=teal!4,colframe=teal!35!black,title={Before you move on}]
What does \mr{तडफडणे} mean? (\textquotedblleft{}To writhe\textquotedblright{}.) Say it once more.
\end{tcolorbox}
\section[\texorpdfstring{\mr{कळवळणे} (kaḷvaḷṇe)}{कळवळणे (kaḷvaḷṇe)}]{\mr{कळवळणे} (kaḷvaḷṇe) — to ache with pity}

\label{lesson:MR-C467-kalvalne}



\begin{quote}
Before the new one: say the Marathi for torment, then the Marathi for to writhe.
\end{quote}

\subsection*{You'll want to know: \mr{कळवळणे}}
\textbf{\mr{कळवळणे}} — \emph{kaḷvaḷṇe} — \textquotedblleft{}to ache with pity\textquotedblright{}.

\begin{quote}
\mr{भुकेल्या} \mr{मुलाला} \mr{पाहून} \mr{आईचा} \mr{जीव} \mr{कळवळला}.
\end{quote}

\begin{quote}
\emph{bhukelyā mulālā pāhūn āīcā jīv kaḷvaḷlā.}
\end{quote}

\begin{quote}
Seeing the hungry child, the mother's heart was wrung.
\end{quote}

Past (he …): \mr{तो} \mr{कळवळला}, \emph{to kaḷvaḷlā}.

Usually of the \mr{जीव} or the \mr{मन}: pity that hurts.

\begin{cousinweb}
\textbf{\mr{कळवळणे}} (\emph{kaḷvaḷṇe}) is a Marathi verb built on a doubled syllable; its origin is not settled.
\end{cousinweb}

\begin{grammarlens}[title={What you've built}]
One more word for describing what you lived and felt.
\end{grammarlens}

\subsection*{Your turn}
\begin{itemize}
\raggedright
  \item \emph{Say it:} \emph{kaḷvaḷṇe}
  \item \emph{Say it:} \emph{kaḷvaḷṇe}, once more
  \item \emph{Say it:} \emph{bhukelyā mulālā pāhūn āīcā jīv kaḷvaḷlā.} — Seeing the hungry child, the mother's heart was wrung.
  \item \emph{Recall:} say the Marathi for torment, then the Marathi for to writhe, then say \emph{kaḷvaḷṇe} again
\end{itemize}

\begin{tcolorbox}[breakable,colback=teal!4,colframe=teal!35!black,title={Before you move on}]
What does \mr{कळवळणे} mean? (\textquotedblleft{}To ache with pity\textquotedblright{}.) Say it once more.
\end{tcolorbox}
\section[\texorpdfstring{\mr{जगणे} (jagṇe)}{जगणे (jagṇe)}]{\mr{जगणे} (jagṇe) — to stay alive; to live}

\label{lesson:MR-C467-jagne-live}



\begin{quote}
Before the new one: say the Marathi for to writhe, then the Marathi for to ache with pity.
\end{quote}

\subsection*{You'll want to know: \mr{जगणे}}
\textbf{\mr{जगणे}} — \emph{jagṇe} — \textquotedblleft{}to stay alive; to live\textquotedblright{}.

\begin{quote}
\mr{पाण्याशिवाय} \mr{कोणीही} \mr{जगू} \mr{शकत} \mr{नाही}.
\end{quote}

\begin{quote}
\emph{pāṇyāśivāy koṇīhī jagū śakat nāhī.}
\end{quote}

\begin{quote}
No one can live without water.
\end{quote}

Past (he …): \mr{तो} \mr{जगला}, \emph{to jaglā}.

Staying alive, where \mr{राहणे} is living somewhere.

\begin{cousinweb}
\textbf{\mr{जगणे}} (\emph{jagṇe}) is a Marathi verb whose origin is not settled; CDIAL lists the older Marathi verb of living, \emph{jiṇẽ}, under Sanskrit \emph{jīvati} (CDIAL 5241).
\end{cousinweb}

\begin{grammarlens}[title={What you've built}]
One more word for describing what you lived and felt.
\end{grammarlens}

\subsection*{Your turn}
\begin{itemize}
\raggedright
  \item \emph{Say it:} \emph{jagṇe}
  \item \emph{Say it:} \emph{jagṇe}, once more
  \item \emph{Say it:} \emph{pāṇyāśivāy koṇīhī jagū śakat nāhī.} — No one can live without water.
  \item \emph{Recall:} say the Marathi for to writhe, then the Marathi for to ache with pity, then say \emph{jagṇe} again
\end{itemize}

\begin{tcolorbox}[breakable,colback=teal!4,colframe=teal!35!black,title={Before you move on}]
What does \mr{जगणे} mean? (\textquotedblleft{}To stay alive; to live\textquotedblright{}.) Say it once more.
\end{tcolorbox}
\section[\texorpdfstring{\mr{डगमगणे} (ḍagmagṇe)}{डगमगणे (ḍagmagṇe)}]{\mr{डगमगणे} (ḍagmagṇe) — to wobble; to waver}

\label{lesson:MR-C467-dagmagne}



\begin{quote}
Before the new one: say the Marathi for to ache with pity, then the Marathi for to stay alive.
\end{quote}

\subsection*{You'll want to know: \mr{डगमगणे}}
\textbf{\mr{डगमगणे}} — \emph{ḍagmagṇe} — \textquotedblleft{}to wobble; to waver\textquotedblright{}.

\begin{quote}
\mr{संकटातही} \mr{तिचा} \mr{धीर} \mr{डगमगला} \mr{नाही}.
\end{quote}

\begin{quote}
\emph{saṅkaṭātahī ticā dhīr ḍagmaglā nāhī.}
\end{quote}

\begin{quote}
Even in the crisis her courage did not waver.
\end{quote}

Past (he …): \mr{तो} \mr{डगमगला}, \emph{to ḍagmaglā}.

Of a boat, a step or resolve: \mr{पाय} \mr{डगमगणे}, to stagger.

\begin{cousinweb}
\textbf{\mr{डगमगणे}} (\emph{ḍagmagṇe}) is inherited: CDIAL lists Marathi \emph{ḍagmagṇẽ}, \textquotedblleft{}to tremble\textquotedblright{}, under the expressive \emph{*ḍag} (CDIAL 5522).
\end{cousinweb}

\begin{grammarlens}[title={What you've built}]
One more word for describing what you lived and felt.
\end{grammarlens}

\subsection*{Your turn}
\begin{itemize}
\raggedright
  \item \emph{Say it:} \emph{ḍagmagṇe}
  \item \emph{Say it:} \emph{ḍagmagṇe}, once more
  \item \emph{Say it:} \emph{saṅkaṭātahī ticā dhīr ḍagmaglā nāhī.} — Even in the crisis her courage did not waver.
  \item \emph{Recall:} say the Marathi for to ache with pity, then the Marathi for to stay alive, then say \emph{ḍagmagṇe} again
\end{itemize}

\begin{tcolorbox}[breakable,colback=teal!4,colframe=teal!35!black,title={Before you move on}]
What does \mr{डगमगणे} mean? (\textquotedblleft{}To wobble; to waver\textquotedblright{}.) Say it once more.
\end{tcolorbox}
